% ============================================================
\section{循環 OLC グラフに対する QUBO アルゴリズム}
% ============================================================

第3章では、候補 overlap 辺を二値変数とする weighted Hamiltonian path の QUBO を用い、
DAG 上での挙動を調べた。
DAG では directed cycle が存在しないため、各 read の入次数と出次数、および選択辺数を
適切に制約すれば、選択辺は一本の Hamilton path として連結される。
一方、実データから得られる OLC グラフには repeat や競合する overlap によって cycle が現れる。
この場合、同じ局所制約を満たしていても、一本の path と複数の cycle が共存する解が残る。
本章では、このような循環グラフを対象として、cycle を排除しながら高重みの Hamilton path を
探索するための QUBO アルゴリズムを構成する。

% ============================================================
\subsection{循環グラフで生じる問題と評価データ}
% ============================================================

候補 overlap 辺 $e=(u,v)\in E$ を使用するかどうかを
$x_e\in\{0,1\}$ で表す。
第3章で用いた edge-variable QUBO の基本形は

\begin{equation}
H_0
=
H_{\mathrm{weight}}
+
H_{\mathrm{degree}}
+
H_{\mathrm{count}}
\label{eq:cyclic_base_hamiltonian}
\end{equation}

である。
$H_{\mathrm{weight}}$ は overlap weight の高い辺を選択する目的項、
$H_{\mathrm{degree}}$ は一つの read に複数の predecessor または successor が接続する状態を抑える項である。
Hamilton path が $N$ 頂点に対して $N-1$ 本の辺を持つことから、

\begin{equation}
H_{\mathrm{count}}
=
A_{\mathrm{count}}
\left(
\sum_{e\in E}x_e-(N-1)
\right)^2
\label{eq:cyclic_count_constraint}
\end{equation}

を加える。

$d_v^{\mathrm{in}}\leq 1$、$d_v^{\mathrm{out}}\leq 1$ が全頂点で成り立つ selected graph は、
directed path と directed cycle の直和として表される。
path component の数を $p$ とすると、path は頂点数より一本少ない辺を持ち、cycle は頂点数と同数の辺を持つため、
selected edge の総数は $N-p$ である。
したがって $\sum_e x_e=N-1$ ならば $p=1$ であるが、cycle の個数は 0 に限定されない。
すなわち式~\eqref{eq:cyclic_base_hamiltonian} の局所制約だけでは、

\begin{equation}
\text{one open path}+\text{one or more directed cycles}
\label{eq:path_plus_cycles}
\end{equation}

を Hamilton path から区別できない。
DAG では存在しなかったこの自由度が、循環 OLC グラフを扱う際の追加条件となる。

本章の評価には CHM13 complex144 を用いる。
これは CHM13 v1.1 の HiFi primary-alignment data から構成した read-level OLC benchmark である。
元データから品質条件を満たす read を抽出して代表骨格を作成し、exact all-vs-all overlap graph の
トポロジーに基づいて複雑な連続区間を選択した後、98\% identity 以上で他の read に完全包含される
9 reads を除去している。
最終的な benchmark は 144 reads、261 directed candidate edges からなり、
最大の nontrivial strongly connected component は 112 頂点である。
全 Hamilton paths を厳密列挙すると 192 通り存在し、単一の単純 cycle ではなく、
多数の分岐と競合する overlap を含む。

benchmark の整理と read orientation の統一には CHM13 reference alignment を利用しているが、
QUBO の edge weight は read--read overlap の alignment のみから計算する。
主結果で用いる edge weight は、matching bases を $M$、alignment block length を $B$、
$d=B-M$ として

\begin{equation}
w_e=M-49d
\label{eq:complex144_weight}
\end{equation}

とした整数値であり、reference coordinate を目的関数へ加えていない。
このため、known reference は benchmark の整理と事後検証に用い、layout score 自体は overlap 情報から定まる。

\begin{figure}[tbp]
\centering
\fbox{\parbox[c][48mm][c]{0.88\linewidth}{\centering
\textbf{図挿入予定：CHM13 complex144 の directed overlap graph}\\[2mm]
144 reads、261 candidate edges、主要 SCC と分岐構造を表示する。}}
\caption{CHM13 complex144 のトポロジー。図は後に差し替える。}
\label{fig:complex144_topology_placeholder}
\end{figure}

% ============================================================
\subsection{順序変数を用いた QUBO}
% ============================================================

cycle を QUBO 内で直接排除する方法として、頂点または辺の順序を二値変数で表す方法がある。
N{\"u}{\ss}lein \textit{et al.}\ \cite{nusslein2022qubo} は Hamiltonian cycle に対して、
選択辺に cycle 上の位置を対応させる order-variable formulation を提案している。
候補辺 $(a,b)$ の order を $P_{ab}$ とすると、$P_{ab}=0$ は未選択、
$P_{ab}>0$ は選択された辺の cycle 上の位置を表す。
これを

\begin{equation}
P_{ab}=\sum_{k=0}^{K-1}2^k p_{ab,k},
\qquad
K=\left\lceil\log_2(N+1)\right\rceil
\label{eq:nusslein_binary_order}
\end{equation}

のように二進展開し、接続する選択辺の order が連続するように制約することで subtour を排除する。
この考え方は Hamiltonian path にも拡張できるが、weighted path では edge selection、始点・終点、
順序 certificate、overlap objective を同時に整合させる必要がある。

本研究では、この order-variable の考え方を残しつつ、Hamilton path の無閉路性に必要な条件を弱めた。
選択された edge $u\rightarrow v$ に対して、rank が厳密に増加すれば directed cycle は存在できるため、

\begin{equation}
x_{uv}=1
\quad\Longrightarrow\quad
P_v>P_u
\label{eq:strict_rank_condition}
\end{equation}

のみを課す。
$P_v=P_u+1$ のように位置を連続させる必要はなく、rank 間に gap があってもよい。
この変更により、rank は経路の絶対位置ではなく、選択辺が一方向に進むことを証明する certificate として扱われる。

各 read $v$ に $K=\lceil\log_2 N\rceil$ 個の rank bits $p_{v,k}$ を割り当て、

\begin{equation}
P_v=\sum_{k=0}^{K-1}2^k p_{v,k}
\label{eq:latent_vertex_rank}
\end{equation}

とする。
式~\eqref{eq:strict_rank_condition} は、$P_v-P_u-1$ の二進減算における borrow を用いて判定する。
edge $e=(u,v)$ の bit $k$ に対する borrow-in を $b$、borrow-out を $c$ とすると、
$c$ は $p_{u,k}$、$b$、$1-p_{v,k}$ の majority である。
この関係は binary variables に対して

\begin{align}
R(u,v,b,c)
={}&u+b+c+ub-uv-bv \notag\\
&-2uc-2bc+2vc
\label{eq:borrow_quadratic_penalty}
\end{align}

という二次式で厳密に表現でき、正しい borrow のときに限り $R=0$ となる。
最下位 bit の borrow-in を 1 として各 bit の borrow を順に計算すると、最終 borrow は

\begin{equation}
b_{e,K}=1
\quad\Longleftrightarrow\quad
P_v\leq P_u
\label{eq:final_borrow_order}
\end{equation}

となる。
したがって selected edge に対してのみ

\begin{equation}
H_{\mathrm{gate}}
=
A_{\mathrm{gate}}
\sum_{e=(u,v)\in E}
 x_e b_{e,K}
\label{eq:strict_rank_gate}
\end{equation}

を課せば、式~\eqref{eq:strict_rank_condition} が成立する。

Hamilton path の始点と終点は固定しない。
各頂点 $v$ に source variable $s_v$ と sink variable $t_v$ を置き、

\begin{align}
s_v+\sum_{u:(u,v)\in E}x_{uv}&=1,\\
t_v+\sum_{w:(v,w)\in E}x_{vw}&=1,
\label{eq:strict_rank_path_degree}
\end{align}

および

\begin{equation}
\sum_v s_v=1,
\qquad
\sum_v t_v=1
\label{eq:strict_rank_endpoint_count}
\end{equation}

を課す。
source と sink が同一頂点になる退化解は $s_vt_v$ に対する penalty で除外する。
これらの次数制約と式~\eqref{eq:strict_rank_condition} を同時に満たす selected graph は、
一つの source から一つの sink まで全頂点を通る Hamilton path となる。
実際、directed cycle が存在すると

\begin{equation}
P_{v_1}<P_{v_2}<\cdots<P_{v_L}<P_{v_1}
\end{equation}

が必要となり、矛盾するためである。
weighted objective は第3章と同様に selected edges の overlap weight を用いる。

% ============================================================
\subsection{complex144 における順序変数 QUBO の評価}
% ============================================================

complex144 では $N=144$、$M=261$ であり、rank bits は一頂点あたり $K=8$ となる。
上記の strict-rank comparator は、261 個の edge variables、288 個の source/sink variables、
1152 個の vertex-rank bits、2088 個の borrow bits を用いるため、論理変数数は

\begin{equation}
Q
=M+2N+NK+MK
=3789
\label{eq:strict_rank_variable_count}
\end{equation}

となる。

この formulation では、order difference を一つの整数平方として展開せず、
bit ごとの borrow relation を局所的な二次式として表す。
そのため、初期に用いた edge-position formulation と比較して、二次係数の最大絶対値を大幅に抑えられた。
表~\ref{tab:order_qubo_comparison} に complex144 での比較を示す。

\begin{table}[tbp]
\centering
\caption{complex144 に対する order-variable QUBO の比較。}
\label{tab:order_qubo_comparison}
\begin{tabular}{lrrr}
\hline
formulation & logical variables & quadratic terms & max $|J|$ \\
\hline
edge-position & 2637 & 77843 & 9437184 \\
strict-rank comparator & 3789 & 33376 & 576 \\
\hline
\end{tabular}
\end{table}

既知の Hamilton path を binary state に符号化した audit では、degree、endpoint、borrow comparator、
selected-edge order のすべてで constraint residual は 0 となった。
また、rank を $0,1,2,\ldots$ と連続に与えた場合だけでなく、全 rank を一定量だけ平行移動した場合や、
順序を保ったまま非一様な gap を入れた場合も同じ objective energy を与えた。
したがって、この formulation は rank の絶対値や間隔ではなく、selected edges 上の順序だけを符号化している。

一方、汎用 QUBO solver を用いた初期テストでは、borrow recurrence 自体は満たせても、
Hamilton path と globally consistent な rank assignment を同時に得ることが難しかった。
基準 penalty では 135--138 本程度の selected edges が得られたが、degree residual と selected-edge order violation が残った。
order gate を強くすると order violation は除けるものの、selected edges が約 80--83 本まで減少し、
一つの Hamilton path へ連結しなかった。
この結果は comparator circuit の不正確さではなく、edge topology を変更すると同時に、
多数の rank bits と borrow bits を整合させなければならないことに起因する。

係数幅を $10^7$ 程度から $10^3$ 未満まで抑えても、logical variables は 3789 個残る。
さらに、二つの Hamilton paths がグラフ上では数本の edge の付け替えで近接していても、
order-variable representation では edge variables に加えて rank certificate も更新する必要がある。
局所的な single-bit update を用いる solver では、この coordinated update が探索障壁となる。
このため、無閉路性を完全に QUBO 内部へ保持する代わりに、
cycle の判定そのものを古典計算へ戻し、QUBO には edge variables だけを残す方法を検討した。

% ============================================================
\subsection{edge-only QUBO と古典計算による cycle feedback}
% ============================================================

edge-only 表現では、探索変数を 261 個の candidate-edge variables $x_e$ に戻し、
式~\eqref{eq:cyclic_base_hamiltonian} を基本 Hamiltonian とする。
この Hamiltonian は cycle を完全には禁止しないが、solver が返した selected graph から
cycle を検出すること自体は難しい最適化問題ではない。
strongly connected components や directed cycle は $O(|V|+|E|)$ のグラフ探索で求められるため、
3789-variable QUBO の内部で rank certificate を毎回更新する必要はない。

そこで、QUBO optimization と classical graph analysis を交互に行う。
第 $t$ 反復で得られた selected graph を $G_x^{(t)}$ とし、そこから検出した directed cycles の集合を
$\mathcal{C}^{(t)}$ とする。
一つの cycle $C$ に対して、特定の一辺を cycle-breaking edge として固定すると、
その選択に依存した bias が生じる。
本研究では、cycle 内の全辺に同じ重みで弱い penalty を分配し、

\begin{equation}
H_{\mathrm{cyc}}(C)
=
\frac{A_{\mathrm{cyc}}}{|C|}
\sum_{e\in C}x_e
\label{eq:uniform_cycle_penalty}
\end{equation}

とする。
これは、長さ $|C|$ の directed cycle に対して可能な $|C|$ 個の cycle-breaking 位置を等価に扱い、
一つの backward edge をあらかじめ指定しない penalty である。
補助変数を追加しないため、cycle feedback を加えても logical variable 数は 261 のままである。

次反復の Hamiltonian は

\begin{equation}
H^{(t+1)}
=
H_0
+
\sum_{C\in\mathcal{C}^{(t)}}
H_{\mathrm{cyc}}(C)
\label{eq:iterative_cycle_hamiltonian}
\end{equation}

とする。
ここで過去に検出した cycle penalty を蓄積するのではなく、各反復で現在の selected graph を解析し、
その時点の cycle に対して Hamiltonian を再構成する。
従って、global order を多数の二値変数として保持する代わりに、
solver が生成した候補解に必要な topology information だけを古典計算で求め、次の QUBO へ戻す構成となる。

\begin{figure}[tbp]
\centering
\begin{tikzpicture}[
  >=Latex,
  box/.style={draw,rounded corners=1.5mm,minimum width=32mm,minimum height=10mm,align=center,font=\small},
  node distance=12mm and 15mm
]
\node[box] (qubo) {QUBO optimization\\$H^{(t)}$};
\node[box,right=of qubo] (state) {selected graph\\$G_x^{(t)}$};
\node[box,right=of state] (check) {classical\\cycle detection};
\node[box,below=of state] (update) {Hamiltonian update\\$H^{(t+1)}$};
\draw[->,thick] (qubo) -- (state);
\draw[->,thick] (state) -- (check);
\draw[->,thick] (check) |- (update);
\draw[->,thick] (update) -| (qubo);
\end{tikzpicture}
\caption{edge-only QUBO と古典的 cycle 判定を交互に用いる反復計算。}
\label{fig:iterative_cycle_feedback}
\end{figure}

この構成の中心は特定の solver に依存しない。
式~\eqref{eq:iterative_cycle_hamiltonian} が QUBO である限り、SQA、量子アニーリング、
SA、tabu search などを各反復の optimizer として用いることができる。
本報告では、order-variable formulation と同一条件で比較しやすいよう OpenJij の SQA を用いて評価する。

% ============================================================
\subsection{数値評価の設定}
% ============================================================

正式評価では、3789-variable strict-rank comparator と、261-variable iterative edge-QUBO を
complex144 上で比較する。
比較の目的は、無閉路性を rank variables として QUBO 内に保持する方法と、
cycle 判定を古典計算へ移して小さい QUBO を反復更新する方法の違いを確認することである。
fixed $H_0$ に対する追加 baseline は置かず、両 formulation に同程度の SQA 計算量を与える。

solver には OpenJij SQASampler の standard single-spin-flip update を用いる。
両方法とも一回の反復で 8 samples を取得し、各 sample 当たり 1400 sweeps、
$\beta=5$、transverse-field scale $\Gamma=1$、Trotter 数 $P=8$ とする。
一つの run は 16 反復で構成し、初回は全変数 0 から開始する。
各反復では 8 samples のうち QUBO energy が最小の状態を次反復へ引き継ぐ。
strict-rank comparator では Hamiltonian を固定したまま計算を継続し、
iterative edge-QUBO ではその間に selected graph の cycle を古典計算で検出して
式~\eqref{eq:iterative_cycle_hamiltonian} を更新する。
cycle penalty は $A_{\mathrm{cyc}}=4$ とし、formal evaluation の前に固定する。

各方法について独立な 24 runs を実行する。
評価時には sampler が返した全 samples から edge selection を取り出し、
Hamilton path の有無と path weight を formulation の補助変数とは独立に検査する。
各 run について、少なくとも一度 Hamilton path を得たか、得られた Hamilton paths の最大 score、
および certified optimum への到達回数を記録する。
score については 24 runs の median、interquartile range、overall best を用いる。
strict-rank comparator ではこれに加えて comparator/order constraint residual を記録し、
iterative edge-QUBO では各反復で検出された cycle 数を記録する。
実行時間も保存するが、主比較は variable 数と得られる Hamilton path の品質に置く。

% TODO(report-20261005):
% 正式 benchmark 実行後、ここに結果表と score distribution を追加する。
% 表の候補列：method / variables / HP runs / median best score / IQR / overall best / optimum hits / runtime.
% strict-rank comparator の既存 audit 値と新しい matched-budget SQA 結果を混同しないこと。

以上の構成により、循環グラフの無閉路性を一つの大きな QUBO に完全符号化する方法と、
小規模な edge-QUBO と古典的な cycle 判定を交互に用いる方法を、同一の実データ上で比較できる。
次章では、この cyclic optimization problem を大規模な OLC graph からどのように切り出すかを扱う。
